\chapter{Ausgangslage}
\label{ch:ausgangslage}

Grundlage dieser Arbeit sind zwei Masterarbeiten, die 2023 an der Hochschule Bochum
parallel entstanden sind: die Hardware des Frequenzumrichters \cite{darici2023} und die
modellbasiert entwickelte Firmware \cite{akguel2023}. Dieses Kapitel beschreibt den
vorgefundenen Stand aus Sicht der Software. Es stellt zunächst das Zielsystem vor, soweit
es für das Verständnis der Firmware nötig ist, dann das Simulink-Modell der Masterarbeit
und schließlich die Befunde, die sich bei der Analyse dieses Modells ergeben haben. Diese
Befunde bestimmen die Reihenfolge der späteren Arbeitsschritte.

% ---------------------------------------------------------------
\section{Zielsystem im Überblick}
\label{sec:zielsystem}

Der Frequenzumrichter besteht aus drei Baugruppen (\cref{fig:zielsystem}):

\begingroup\sloppy
\begin{description}
  \item[Steuerplatine] Eigenentwicklung aus \cite{darici2023} mit einem
    \textbf{TMS320F28386D} \cite{ti_f2838x_dm}. Der Baustein
    enthält zwei C28x-Kerne (CPU1, CPU2) mit je \SI{200}{\mega\hertz} und einen
    Arm-Cortex-M4-Kern, den \ac{CM}, für Ethernet, der über \ac{IPC} Daten mit den C28x-Kernen
    austauscht. Die Firmware läuft ausschließlich auf CPU1.
    Den Systemtakt erzeugt ein externer \SI{25}{\mega\hertz}-Oszillator. Über USB-C stellt
    ein FT2232HL eine \acs{JTAG}-Debugschnittstelle und eine serielle
    Schnittstelle bereit, über einen galvanisch getrennten Transceiver ist \ac{CAN}
    herausgeführt.
  \item[Leistungsstufe] Das zugekaufte \ac{EVM} \textbf{DRV8300DIPW-EVM}
    \cite{ti_slvubv6} mit dreiphasiger MOSFET-Brücke, Gatetreiber DRV8300, drei
    Low-Side-Shunts mit \ac{CSA}, Temperatursensor und dem Anschluss für die
    Hall-Sensoren. Es ist über zwei Stiftleisten im Raster des LaunchPads F280049C mit
    der Steuerplatine verbunden.
  \item[Motor] Der Motor ist ein permanentmagneterregter Synchronmotor vom Typ \textbf{ebm-papst ECI~63.60} \cite{ebmpapst_eci6360}. Er ist
     für \SI{24}{\volt} ausgelegt und besitzt vier Polpaare. Drei digitale Hall-Sensoren ermöglichen eine Lageerfassung der Ausgangswelle.
\end{description}
\endgroup

\begin{figure}[htbp]
  \centering
  \begin{tikzpicture}[
      node distance=8mm and 14mm,
      box/.style={draw, rounded corners=2pt, align=center, minimum height=26mm,
                  text width=34mm, font=\small},
      sig/.style={-Latex, font=\scriptsize},
    ]
    \node[box] (mcu) {\textbf{Steuerplatine}\\TMS320F28386D\\CPU1: Regelung\\
      \SI{25}{\mega\hertz} ext.\ Oszillator};
    \node[box, right=22mm of mcu] (evm) {\textbf{DRV8300DIPW-EVM}\\Gatetreiber\\
      B6-Brücke\\3 Shunts + CSA};
    \node[box, right=14mm of evm, text width=22mm] (mot) {\textbf{ECI 63.60}\\PMSM\\
      4 Polpaare\\Hall A/B/C};
    \node[box, below=10mm of mcu, minimum height=9mm] (pc) {PC / Bedienung};
    \node[box, below=10mm of evm, minimum height=9mm] (psu) {Labornetzteil
      \SI{24}{\volt}};

    \draw[sig] ([yshift=8mm]mcu.east) -- node[above]{6$\times$ PWM} ([yshift=8mm]evm.west);
    \draw[sig] ([yshift=-1mm]evm.west) -- node[above]{$i_a, i_b, i_c$, PVDD}
      ([yshift=-1mm]mcu.east);
    \draw[sig] ([yshift=-9mm]evm.west) -- node[above]{Hall, Temp.} ([yshift=-9mm]mcu.east);
    \draw[sig] ([yshift=4mm]evm.east) -- node[above]{$u_{abc}$} ([yshift=4mm]mot.west);
    \draw[sig] ([yshift=-6mm]mot.west) -- node[below]{Hall} ([yshift=-6mm]evm.east);
    \draw[Latex-Latex, font=\scriptsize] (pc) -- node[right]{CAN, USB (JTAG/UART)} (mcu);
    \draw[sig] (psu) -- node[right]{PVDD} (evm);
  \end{tikzpicture}
  \caption{Aufbau des Zielsystems mit den für die Firmware relevanten Signalen}
  \label{fig:zielsystem}
\end{figure}

Die für die Firmware wichtigen Peripherien sind in \cref{tab:peripherie} zusammengefasst.
Die drei Brückenzweige werden von ePWM6 bis ePWM8 angesteuert, mit \SI{20}{\kilo\hertz}
im Up-Down-Zählbetrieb (mittensymmetrische \ac{PWM}). Die Totzeit erzeugt der Gatetreiber,
nicht der Mikrocontroller. Eine Besonderheit des DRV8300\textbf{DI} ist, dass er die
Low-Side-Eingänge INLx invertiert. Der Mikrocontroller gibt deshalb auf ePWMxA und ePWMxB
dasselbe Signal aus. ePWM6 startet zugleich die Wandlung der Phasenströme, und zwar dann,
wenn die Low-Side-Schalter leiten und die Shunts den Strom sehen. Der Interrupt am Ende
dieser Wandlung ruft die Stromregelung auf. Damit ist die Stromregelung fest an die
\ac{PWM}-Periode von \SI{50}{\micro\second} gekoppelt.

\begin{table}[htbp]
  \centering
  \caption{Von der Firmware der \ac{MA} genutzte Peripherie (CPU1)}
  \label{tab:peripherie}
  \small
  \begin{tabularx}{\linewidth}{@{}>{\raggedright}p{36mm}>{\raggedright}p{34mm}
      >{\raggedright\arraybackslash}X@{}}
    \toprule
    Funktion & Peripherie & Bemerkung \\
    \midrule
    Brückenansteuerung & ePWM6/7/8 (GPIO10--15) & \SI{20}{\kilo\hertz}, Periode 5000 Takte \\
    Phasenströme $i_a$, $i_b$ & ADCINB0, ADCINB2 & Trigger ePWM6-SOCA, Interrupt ADCB1 $\to$ Stromregelung \\
    Phasenstrom $i_c$ & ADCINB3 & verdrahtet, aber nicht genutzt \\
    Zwischenkreisspannung, GVDD & ADCINA5, ADCINC3 & Unterspannungserkennung \\
    MOSFET-Temperatur & ADCINC2 & TMP23x, Übertemperatur \\
    Drehzahlsollwert & ADCINA0 & Potentiometer \\
    Hall A/B/C & eCAP1--3 (GPIO24--26) & Interrupt bei jedem Flankenwechsel \\
    SafeState-Taster & GPIO30 & Not-Aus in Software \\
    Kommunikation & eCAN-A, IPC zum CM & CAN-Kommandos, Ethernet-Logging über den CM \\
    \bottomrule
  \end{tabularx}
\end{table}

Einige Eigenschaften der Hardware wirken sich unmittelbar auf die Software aus. Dazu
zählen die Speisung der ADC-Referenz aus der \SI{3,3}{\volt}-Schaltreglerschiene und die
nur eingeschränkt realisierbare Überstromabschaltung in Hardware. Die Analyse dieser
Eigenschaften ist Gegenstand des Hardwareteils \hwteil{Hardware-Review}. Im vorliegenden
Bericht werden sie nur insoweit betrachtet, als sie Entwurf und Verhalten der Firmware
beeinflussen.

% ---------------------------------------------------------------
\section{Simulink-Modell der Masterarbeit}
\label{sec:impmodel}

Die Firmware der \ac{MA} \cite{akguel2023} enthält keinen handgeschriebenen C-Code. Sie
wird vollständig mit Embedded Coder \cite{mw_ecoder} und dem C2000 Microcontroller Blockset
\cite{mw_c2000} aus dem Modell \datei{impmodel.slx} generiert (MATLAB R2023a). Die Parameter
(Motor, Leistungsstufe, Bezugsgrößen, Reglerverstärkungen, Abtastzeiten) liegen in der
Binärdatei \datei{params\_f2838.mat}. Beide Dateien sind im Repository unverändert unter
\datei{matlab/models/legacy/} abgelegt und dienen in dieser Arbeit als Referenz für alle
Vergleiche. \Cref{fig:impmodel} zeigt die oberste Modellebene.

\begin{figure}[htbp]
  \centering
  \begin{tikzpicture}[
      font=\footnotesize,
      sub/.style={draw, align=center, minimum height=12mm, text width=22mm,
                  font=\footnotesize\sffamily},
      irq/.style={draw, dashed, align=center, minimum height=9mm, text width=22mm,
                  font=\footnotesize\sffamily},
      aux/.style={draw, align=center, minimum height=7mm, text width=22mm,
                  font=\footnotesize\sffamily},
      sig/.style={-Latex},
      lab/.style={font=\scriptsize, align=center},
    ]
    % Hauptsignalfluss
    \node[sub] (in) {Input\\Controller};
    \node[sub, right=16mm of in] (spd) {Speed Control\\\scriptsize $T_s = \SI{0,5}{ms}$};
    \node[sub, right=16mm of spd] (trq) {Torque Controller\\\scriptsize $T_s = \SI{50}{\micro s}$};
    \node[sub, right=16mm of trq] (plant) {Inverter and\\Motor -- Plant\\Model};

    % Interruptgesteuerte Subsysteme
    \node[irq, above=14mm of trq] (cg) {Code\\Generation};
    \node[irq, above=14mm of spd] (hall) {Hall Sensor\\A / B / C};
    \node[irq, above=14mm of plant] (ipc) {IPC Write for\\eCAP INTs};

    % Initialisierung und Hilfsfunktionen
    \node[aux, below=12mm of in] (init) {System Initialize};
    \node[aux, below=12mm of spd] (off) {ADC Offset Function};
    \node[aux, below=12mm of plant] (hb) {Heartbeat LED};

    \draw[sig] (in) -- node[lab, above]{Sollwert,} node[lab, below]{Freemode} (spd);
    \draw[sig] (spd) -- node[lab, above]{$i_{dq,\mathrm{ref}}$,} node[lab, below]{Freemode} (trq);
    \draw[sig] ([yshift=3mm]trq.east) -- node[lab, above]{Tastgrade} ([yshift=3mm]plant.west);
    \draw[sig] ([yshift=-3mm]plant.west) -- node[lab, below]{$i_{ab}$, Lage} ([yshift=-3mm]trq.east);
    \draw[sig] (trq.south) -- ++(0,-5mm) -| node[lab, pos=0.25, below]{Drehzahl} (spd.south);
    \draw[sig] (cg) -- node[lab, right]{ADCB1} (trq);
    \draw[sig] (cg) -- node[lab, above]{eCAP1--3} (hall);
    \draw[sig, densely dotted] (hall) -- node[lab, right]{Data Stores} (spd);
  \end{tikzpicture}
  \caption{Vereinfachte oberste Ebene des Modells \datei{impmodel.slx} aus \cite{akguel2023}.
    Gestrichelt: durch Interrupts aufgerufene Subsysteme. Die Original-Modellansicht zeigt
    \cref{fig:impmodel_orig} im Anhang.}
  \label{fig:impmodel}
\end{figure}

\subsection{Struktur}

Das Modell gliedert sich in folgende Subsysteme:

\begingroup\sloppy
\begin{description}
  \item[\block{System Initialize}, \block{ADC Offset Function}] Initialisierung, unter
    anderem der Pull-up-Widerstände über direkte Registerzugriffe, und Kalibrierung des
    Stromsensor-Offsets aus acht Abtastwerten mit Plausibilitätsprüfung.
  \item[\block{Code Generation}] Die Hardware-Interrupts: ADCB1 ruft den
    \block{Torque Controller} auf, eCAP1 bis eCAP3 die Hall-Auswertung.
  \item[\block{Hall Sensor A/B/C}] Interrupt-Routinen je Hall-Signal. Sie lesen den
    Hall-Zustand, prüfen seine Gültigkeit und legen Zustand, Drehrichtung und Zeitstempel
    in globalen Data Stores ab.
  \item[\block{Input Controller}] Sollwert vom Potentiometer, Unterspannung,
    Übertemperatur, SafeState-Taster sowie die CAN-Kommandos \enquote{Freemode} und
    \enquote{speedMode}. Außerdem werden Zwischenkreis- und Gatetreiberspannung zyklisch
    über CAN gesendet.
  \item[\block{Speed Control}] Drehzahlregelung, Abtastzeit \SI{0,5}{\milli\second}. Sie
    enthält die Freemode-Logik (\block{Digital\_Freemode\_Check} und
    \block{Mechanical\_Freemode\_Check}), die Sollwertberechnung (\block{Speed calculator})
    und einen \ac{PI} als Drehzahlregler (\block{PI\_Controller\_Speed}). Ausgang ist der
    Stromsollwert $\boldsymbol{i}_{dq,\mathrm{ref}} = (0,\ i_{q,\mathrm{ref}})$.
  \item[\block{Torque Controller}] Die eigentliche \ac{FOC}. Sie wird alle
    \SI{50}{\micro\second} von der \ac{ISR} des \ac{ADC} aufgerufen und umfasst
    Stromerfassung und Skalierung, Lage und Drehzahl aus den Hall-Signalen, Clarke- und
    Park-Transformation, \acs{PI}-Stromregler für $d$ und $q$,
    Spannungsbegrenzung (\enquote{DQ-Limiter}), inverse Park-Transformation,
    Raumzeigermodulation und Umrechnung in Vergleichswerte der ePWM-Module.
  \item[\block{Inverter and Motor -- Plant Model}] Streckenmodell aus Wechselrichter
    (Mittelwertmodell) und Motor für die Simulation.
  \item[\block{IPC Write for eCAP INTs}, \block{Heartbeat LED}] Übertragung von Messwerten
    an den \ac{CM}-Kern, der sie über ein zweites Modell (\datei{TCP\_Server.slx}) per
    Ethernet an den PC sendet, sowie eine Blink-LED als Lebenszeichen.
\end{description}
\endgroup

Die Regelung ist eine klassische Kaskade aus äußerem Drehzahl- und innerem Stromregelkreis
ohne Feldschwächung ($i_{d,\mathrm{ref}} = 0$). Gerechnet wird in normierten Größen, angegeben in
\ac{PU}, überwiegend im Datentyp \texttt{single} oder \texttt{fixdt(1,32,17)}. Die
Reglerverstärkungen wurden nicht von Hand eingestellt, sondern mit der internen, nicht
dokumentierten Funktion \texttt{mcb.internal.SetControllerParameters} des Motor Control
Blockset aus den Motordaten berechnet.

Alle hardwarenahen Teile sind über Variantenblöcke doppelt ausgeführt: Die Variante
\texttt{sim} speist die Regelung aus dem Streckenmodell, die Variante \texttt{codegen}
aus den Treiberblöcken des C2000 Blocksets. Simulations- und Codegenerierungspfad teilen
sich also die Reglerlogik, nicht aber die Messkette. Diese Trennung wird bei der Diagnose
in \cref{ch:regler} eine Rolle spielen.

\subsection{Sicherheitslogik}

Die einzige Rückfallebene ist der \enquote{Freemode}: Alle Schalter der Brücke werden
geöffnet, der Motor läuft frei aus, und der Fehler des Drehzahlreglers wird auf null
gesetzt. Ausgelöst wird er durch Unterspannung, Übertemperatur (bis zum Neustart
gespeichert), ein CAN-Kommando oder den SafeState-Taster. Eine Überspannungs- oder
Überstromabschaltung gibt es nicht. Die Untergrenze des Drehzahlreglers ist fest auf
$-0{,}5$ gesetzt, was ein leichtes Bremsen erlaubt.

\subsection{Anforderungen und Tests}

Die \ac{MA} leitet aus einem Lastenheft neun Systemmodellanforderungen
(SW\_ANF\_01 bis SW\_ANF\_09) ab. Vier davon sind mit Simulink Test als \ac{MIL}-Tests an
Test-Harnesses im Modell abgesichert (\cref{tab:mil_ma}), alle vier bestanden. Die
feldorientierte Regelung selbst (SW\_ANF\_06) und die Hall-Auswertung (SW\_ANF\_09) sind
nur auf der Hardware gezeigt, nicht modellbasiert getestet. Die Testdatei
\datei{Testfile\_F2838.mldatx} ist Teil der \ac{MA}.

\begin{table}[htbp]
  \centering
  \caption{\ac{MIL}-Tests der \ac{MA} \cite{akguel2023}}
  \label{tab:mil_ma}
  \small
  \begin{tabularx}{\linewidth}{@{}llX@{}}
    \toprule
    Harness & Getesteter Block & Anforderung \\
    \midrule
    \texttt{Freemode\_Check\_TestHarness} & \block{Digital\_Freemode\_Check}
      & SW\_ANF\_01/07/08 (Freemode bei CAN, Unterspannung, Übertemperatur) \\
    \texttt{SafeState\_TestHarness} & \block{Mechanical\_Freemode\_Check}
      & SW\_ANF\_03 (SafeState-Taster) \\
    \texttt{MaxVel\_TestHarness} & \block{Speed calculator}
      & SW\_ANF\_04 (Höchstdrehzahl je speedMode) \\
    \texttt{SpeedController\_TestHarness} & \block{PI\_Controller\_Speed}
      & SW\_ANF\_05 (Drehzahlregler, Begrenzung) \\
    \bottomrule
  \end{tabularx}
\end{table}

\subsection{Was nicht übernommen werden konnte}

Zwei Teile des Ausgangsstands ließen sich nicht direkt weiterverwenden:

\begin{itemize}
  \item \textbf{Ethernet-Logging:} Weil die Ethernet-Pins auf der eigenen Platine anders
    belegt sind als auf der controlCARD des Herstellers, wurden in der \ac{MA} Dateien des
    C2000 Support Package geändert (Präprozessorschalter
    \texttt{CARPEDIEM\_PIN\_CORRECTION}). Diese Änderungen sind nicht Teil des Archivs.
    Außerdem verstößt eine Änderung an der Toolbox gegen die Reproduzierbarkeit des Builds.
    Die Ethernet-Strecke über den \ac{CM} wurde daher nicht weiterverfolgt; die
    Kommunikation läuft stattdessen über \ac{CAN} (\cref{ch:kommunikation}).
  \item \textbf{Zielprozessor:} Das Modell ist für den TMS320F28388D der controlCARD
    konfiguriert, auf der Platine sitzt der pinkompatible TMS320F28386D. Auch die
    Taktquelle in den Hardware-Einstellungen musste an den externen
    \SI{25}{\mega\hertz}-Oszillator angepasst werden (\cref{ch:codegen}).
\end{itemize}

% ---------------------------------------------------------------
\section{Befunde aus der Modellanalyse}
\label{sec:befunde}

Bevor das Modell umgebaut wurde, wurden Reglerblöcke, Parameterdatei und Messkette im
Detail untersucht. Die Ergebnisse sind in \cref{tab:befunde} als Befunde B1 bis B6
zusammengefasst. Sie werden im weiteren Bericht mit diesen Kürzeln referenziert.

\begin{table}[htbp]
  \centering
  \caption{Befunde aus der Analyse des Altmodells}
  \label{tab:befunde}
  \small
  \renewcommand{\arraystretch}{1.15}
  \begin{tabularx}{\linewidth}{@{}l>{\raggedright\arraybackslash}X
      >{\raggedright\arraybackslash}X >{\raggedright\arraybackslash}p{27mm}@{}}
    \toprule
    Nr. & Befund & Wirkung & Behandlung \\
    \midrule
    B1 & I-Anteil des Drehzahlreglers um Größenordnungen zu klein
       & praktisch P-Regler, bleibende Drehzahlabweichung
       & \cref{ch:regler} \\
    B2 & Stromregler: I-Anteil wie B1 zu klein, $K_p$ von Hand verändert,
         Ausgang auf $[0, 1]$ begrenzt
       & negative Stellspannung unmöglich
       & \cref{ch:regler} \\
    B3 & Motordaten um Faktoren neben dem Datenblatt
       & Auslegung und Streckenmodell passen nicht zum Motor
       & Profil \texttt{datasheet} (\cref{ch:vorgehen}) \\
    B4 & Strom-PU-Basis im Modell weicht von der Bezugsgröße ab
       & wirksame Reglerverstärkung anders als ausgelegt
       & \cref{ch:regler} \\
    B5 & Variable \texttt{dataType} fehlt in der Parameterdatei
       & Modell nicht ohne Zusatzschritt übersetzbar
       & \cref{ch:bibliothek} \\
    B6 & ADC-Referenz \SI{3,36}{\volt} statt \SI{3,3}{\volt}
       & \SI{1,8}{\percent} Verstärkungsfehler
       & Profil \texttt{datasheet} \\
    \bottomrule
  \end{tabularx}
\end{table}

\subsection*{B1 und B2: Integralanteil der Regler}

Am Block \block{PI\_Controller\_Speed} ist als Integralverstärkung
\texttt{Ki\_speed\_texas*Ts\_speed} eingetragen, gleichzeitig ist die Option
\texttt{UseKiTs} gesetzt. Die Variable \texttt{Ki\_speed\_texas} stammt aus
\texttt{SetControllerParameters} und ist bereits der Integralbeiwert je Abtastschritt in der
Form der TI-Bibliothek, also $T_s/T_n$ ohne $K_p$. Er wird im Block ein zweites Mal mit der
Abtastzeit multipliziert. Statt etwa $0{,}095$ je Schritt wirkt so ein Wert von etwa
$6{,}6 \cdot 10^{-6}$. Der Drehzahlregler arbeitet damit praktisch als reiner P-Regler mit
$K_p = 7{,}2$ und kann eine bleibende Regelabweichung nicht ausregeln.

Die beiden Stromregler im \block{Torque Controller} haben denselben Fehler im
Integralanteil. Dazu kommen zwei weitere Eingriffe: Der Proportionalbeiwert ist als
\texttt{Kp\_i - 0.0002} eingetragen, offenbar eine Handkorrektur, und der Reglerausgang ist
auf $[0, 1]$ begrenzt. Im Betrieb braucht die $d$-Achse aber eine negative Spannung, etwa
$u_d \approx -\omega L i_q$. Die Begrenzung verhindert das.

Ob \texttt{UseKiTs} bedeutet, dass der eingetragene Wert bereits $K_i T_s$ ist, oder dass der
Block intern mit $T_s$ multipliziert, ändert nur die Größe des Fehlers. In beiden Lesarten
ist der Integralanteil um mindestens drei Größenordnungen zu klein.

\subsection*{B3: Motordaten}

Die Motordaten in \datei{params\_f2838.mat} weichen deutlich vom Datenblatt ab
(\cref{tab:motordaten}). Die ältere Parameterdatei der \ac{MA} für das LaunchPad F280049C
enthält dagegen die Datenblattwerte. Aus Spannungskonstante und Leerlaufdrehzahl
(\SI{6000}{\per\minute} bei \SI{24}{\volt}) ergibt sich ein Polradfluss von etwa
\SI{5,9}{\milli\volt\second}. Widerstand und Induktivität im Datenblatt sind sehr
wahrscheinlich Werte zwischen zwei Klemmen, je Strang also die Hälfte. Die
Reglerverstärkungen der \ac{MA} sind aus den abweichenden Werten berechnet. Der
Proportionalbeiwert des Stromreglers ist dadurch etwa um den Faktor 500 kleiner als bei
Auslegung mit den Datenblattwerten. Motorparameter gemessen hat die \ac{MA} nicht: Der Autor
nennt das selbst als Schwachstelle \cite{akguel2023}.

\begin{table}[htbp]
  \centering
  \caption{Motordaten ECI~63.60: Parameterdatei der \ac{MA} und Ableitung aus dem
    Datenblatt \cite{ebmpapst_eci6360}}
  \label{tab:motordaten}
  \small
  \begin{tabular}{@{}lrr@{}}
    \toprule
    Größe & \datei{params\_f2838.mat} & aus Datenblatt, je Strang \\
    \midrule
    Strangwiderstand $R_s$ & \SI{9}{\milli\ohm} & \SI{22,5}{\milli\ohm} \\
    Induktivität $L_d = L_q$ & \SI{0,153}{\micro\henry} & \SI{38}{\micro\henry} \\
    Polradfluss $\Psi_\mathrm{PM}$ & \SI{0,69}{\milli\volt\second} & \SI{5,9}{\milli\volt\second} \\
    \bottomrule
  \end{tabular}
\end{table}

\subsection*{B4: Skalierung der Strommessung}

Die Bezugsgrößen des Modells (Struktur \texttt{PU\_System}) setzen
$I_\mathrm{base} = \SI{17,6}{\ampere}$ und $U_\mathrm{base} = \SI{20,78}{\volt}$ an, und mit
diesen Werten werden die Regler ausgelegt. Im Codegenerierungspfad wird der Strom aber
direkt als $(\text{ADC} - \text{Offset})/2048$ berechnet. Ein \ac{PU} ist dort also der
Strom, bei dem der Strommessverstärker \SI{1,65}{\volt} Hub liefert. \texttt{PU\_System}
kommt in diesem Pfad nicht vor. Der Parameter \texttt{ISenseVoltPerAmp}
(\SI{3}{\milli\volt\per\ampere}) berücksichtigt nur den Shunt, nicht die Verstärkung des
\ac{CSA} auf dem \ac{EVM}. Welcher Strom einem \ac{PU} tatsächlich entspricht und welches
Vorzeichen die Messung hat, war zu Beginn offen. Geklärt wurde das erst mit den Unterlagen
des \ac{EVM} \cite{ti_slvubv6, ti_slvc817} im Zuge der Reglerkorrektur
(\cref{ch:regler}).

\subsection*{B5 und B6: Parameterdatei und ADC-Referenz}

Das Modell verwendet eine Variable \texttt{dataType} zur Wahl zwischen Gleit- und
Festkommavarianten, die in \datei{params\_f2838.mat} fehlt. Ohne zusätzliches Anlegen dieser
Variable lässt sich das Modell nicht übersetzen. Die Referenzspannung der \ac{ADC} ist in den
Parametern mit \SI{3,3}{\volt} angesetzt. Auf der Platine hängt sie an der
\SI{3,3}{\volt}-Versorgung, die mit \SI{3,36}{\volt} gemessen wurde \cite{darici2023}. Alle
analogen Messwerte haben dadurch einen Verstärkungsfehler von \SI{1,8}{\percent}.

\subsection*{Weitere Auffälligkeiten}

Neben den Befunden B1 bis B6 fielen kleinere Unstimmigkeiten auf, die keinen eigenen
Arbeitsschritt erforderten, aber bei einer Weiterentwicklung zu beachten sind:

\begin{itemize}
  \item Der Hall-Lageoffset ist doppelt gepflegt: \num{0,411} in den Motorparametern,
    \num{0,38} fest im Block \block{Mechanical to Electrical Position}. Der Wert im Block
    wurde laut \ac{MA} experimentell ermittelt.
  \item Die Mailbox für das CAN-Kommando \enquote{speedMode} ist in der
    Schnittstellentabelle der \ac{MA} mit 10, im Modell mit 15 angegeben.
  \item Mehrere Anforderungstexte entsprechen nicht mehr dem Modell, zum Beispiel die
    Begrenzung des Drehzahlreglers in SW\_ANF\_05.
  \item Die \ac{IPC}-Blöcke senden ständig an den \ac{CM}-Kern, auch wenn dort kein Programm
    läuft. Das führte später auf der Hardware dazu, dass die Strom-\ac{ISR} hängen blieb
    (\cref{ch:kommunikation}).
\end{itemize}

\section{Folgerungen}
\label{sec:folgerungen}

Der Ausgangsstand ist eine lauffähige, in sich geschlossene Firmware mit einer
nachvollziehbaren Struktur und einem Grundstock an \ac{MIL}-Tests. Für eine
Weiterverwendung als Framework fehlt aber dreierlei: 
\begin{enumerate}
	\item Die Parameter sind nicht
	nachvollziehbar (Binärdatei, interne Auslegungsfunktion, Befunde B3 bis B6).
	\item Die Regler
	arbeiten nicht wie ausgelegt (B1, B2).
	\item Die Komponenten sind fest in ein einziges
	Anwendungsmodell eingebaut.
\end{enumerate}
 Daraus ergibt sich die Reihenfolge der folgenden Kapitel:
zuerst die Parameter aus Skripten bei bitgleichem Verhalten, dann die Aufteilung in eine
Bibliothek, erst danach die Korrekturen an den Reglern. Jede Korrektur wurde einzeln vorgenommen und in anschließend in der Simulation nachgewiesen.
Dadurch war es möglich die Korrektheit jeder Korrektur zu überprüfen bevor die nächste vorgenommen wurde.
